% §2 Background
% Status: Written (stub)

\section{Background}
\label{sec:background}

\subsection{Performance Portability}

Performance portability (PP) quantifies how well a single codebase performs across
different hardware platforms. Following~\cite{Pennycook2019}, we use the
Performance Portability Coefficient (PPC):

\begin{equation}
\label{eq:ppc}
\text{PPC}(a, p) = \frac{|H|}{\sum_{h \in H} \frac{1}{e_h(a, p)}}
\end{equation}

where $H$ is the set of platforms with working implementations, and $e_h(a, p)$
is the efficiency of abstraction $a$ on problem $p$ relative to the native baseline
on platform $h$.

\subsection{GPU Abstraction Layers}

\paragraph{Kokkos~\cite{Carter2014}.}
Data structure abstraction built on portable Views and Execution Spaces.
Widely adopted in ECP codes ($\sim$30\% per Evans et al.~\cite{Evans2022}).

\paragraph{RAJA~\cite{Hornung2014}.}
Loop abstraction: execution policies, index sets, kernel fusion.
Different philosophy from Kokkos --- loop-level vs.\ data-structure-level.

\paragraph{SYCL~\cite{Reyes2020}.}
Open Khronos standard. Vendor-neutral heterogeneous programming.
Implementations: Intel DPC++, AdaptiveCpp (hipSYCL).

\paragraph{Julia~\cite{Bezanson2017}.}
High-level language with JIT compilation. CUDA.jl/AMDGPU.jl backends.
Godoy et al.~\cite{Godoy2023} show Julia achieves $>0.9$ efficiency on AMD
but struggles on NVIDIA for DGEMM --- a pattern we investigate.

\subsection{Hardware Platforms}

TODO: Table of three target platforms (A100, MI250X, PVC) with peak BW and FLOP/s.
See project\_spec.md \S5.1.
